\selectlanguage{american}%

\section{State of the Art}

We have now seen various algorithms to minimize a convex function.They
have different runtimes and seem difficult to compare. Before going
deeper and discussing more results, here we discuss how to understand
these different results and their relationships, via reductions between
them. We will do this at three different levels. The most abstract
is the level of oracles, which we will define presently; the second
is black-box algorithms which access the function in restricted ways,
e.g., via the gradient; and the last is with full information and
under additional structural assumptions.

Here we summarize the state-of-the-art in two tables. We tried to
keep the tables short, and there are many important results omitted.

\begin{table}[H]
\begin{centering}
\begin{tabular}{|>{\centering}m{0.15\paperwidth}|>{\centering}m{0.15\paperwidth}|>{\centering}m{0.15\paperwidth}|>{\centering}m{0.1\paperwidth}|>{\centering}m{0.15\paperwidth}|}
\hline 
{\scriptsize Convex Problem} & {\scriptsize Assumption} & {\scriptsize Algorithm} & {\scriptsize Iterations} & {\scriptsize Cost/Iter}\tabularnewline
\hline 
{\scriptsize
\[
\min_{x}f(x)
\]
} & \centering{}{\scriptsize -} & {\scriptsize Hybrid Center Cutting Plane} & {\scriptsize$n\log(\frac{n}{\epsilon})$} & {\scriptsize$\mathcal{T}_{f}+n^{2}\log^{O(1)}(\frac{n}{\epsilon})$}\tabularnewline
\hline 
{\scriptsize
\[
\min_{x\in\Omega_{1}\cap\Omega_{2}}c^{\top}x
\]
} & \centering{}{\scriptsize -} & {\scriptsize Hybrid Center Cutting Plane} & {\scriptsize$n\log(\frac{n}{\epsilon})$} & {\scriptsize$\mathcal{T}_{\Omega_{1}}+\mathcal{T}_{\Omega_{2}}+n^{2}\log^{O(1)}(\frac{n}{\epsilon})$}\tabularnewline
\hline 
{\scriptsize
\[
\min_{x}f(x)+g(x)
\]
} & {\scriptsize$\|\nabla f\|_{2}\leq G$, $\|x^{*}\|_{2}\leq R$} & {\scriptsize Gradient Descent} & {\scriptsize$\left(\frac{GR}{\epsilon}\right)^{2}$} & {\scriptsize$\mathcal{T}_{\nabla f}+\mathcal{T}_{g}$}\tabularnewline
\hline 
{\scriptsize
\[
\min_{x}f(x)+g(x)
\]
} & {\scriptsize$\|\nabla f\|_{2}\leq G$, $\nabla^{2}f\succeq\mu\cdot I$} & {\scriptsize Gradient Descent} & {\scriptsize$\frac{G^{2}}{\mu\epsilon}$} & {\scriptsize$\mathcal{T}_{\nabla f}+\mathcal{T}_{g}$}\tabularnewline
\hline 
{\scriptsize
\[
\min_{x}f(x)+g(x)
\]
} & \centering{}{\scriptsize$\nabla^{2}f\preceq L\cdot I$, $\|x^{*}\|_{2}\leq R$} & {\scriptsize Accelerated Gradient Descent} & {\scriptsize$\sqrt{\frac{LR^{2}}{\epsilon}}$} & {\scriptsize$\mathcal{T}_{\nabla f}+\mathcal{T}_{g}$}\tabularnewline
\hline 
{\scriptsize
\[
\min_{x}f(x)+g(x)
\]
} & \centering{}{\scriptsize$\mu\cdot I\preceq\nabla^{2}f\preceq L\cdot I$,
$\|x^{*}\|_{2}\leq R$} & {\scriptsize Accelerated Gradient Descent} & {\scriptsize$\sqrt{\frac{L}{\mu}}\log(\frac{LR^{2}}{\epsilon})$} & {\scriptsize$\mathcal{T}_{\nabla f}+\mathcal{T}_{g}$}\tabularnewline
\hline 
{\scriptsize
\[
\min_{x\in\Omega}f(x)
\]
} & {\scriptsize$\nabla f$ is $L$ Lipschitz in $\|\cdot\|$, }{\scriptsize\par}

{\scriptsize$\max_{x,y\in\Omega}\|x-y\|\leq R$} & {\scriptsize Frank Wolfe} & {\scriptsize$\frac{LR^{2}}{\epsilon}$} & {\scriptsize$\mathcal{T}_{\nabla f}+\mathcal{T}_{\Omega}$ }\tabularnewline
\hline 
{\scriptsize
\[
\min_{x}\E_{\omega}f_{\omega}(x)
\]
} & {\scriptsize$\E_{\omega}\|\nabla f_{\omega}\|_{2}^{2}\leq G^{2}$,
$\|x^{*}\|_{2}\leq R$} & {\scriptsize Stochastic Gradient Descent} & {\scriptsize$\left(\frac{GR}{\epsilon}\right)^{2}$} & {\scriptsize$\mathcal{T}_{\nabla f_{\omega}}$}\tabularnewline
\hline 
{\scriptsize
\[
\min_{x}f(x)\defeq\frac{1}{n}\sum_{i}f_{i}(x)
\]
} & {\scriptsize$\nabla^{2}f\succeq\mu\cdot I,\nabla^{2}f_{i}\preceq L_{i}\cdot I$} & {\scriptsize Accelerated Stochastic Descent} & \multicolumn{2}{c|}{{\scriptsize Total Cost: $(\mathcal{T}_{\nabla f}+\sqrt{\frac{1}{\mu}\sum_{i}L_{i}}\mathcal{T}_{\nabla f_{i}})\log(\frac{1}{\epsilon})$}}\tabularnewline
\hline 
{\scriptsize
\[
\min_{x\in\Rn}f(x)
\]
} & {\scriptsize$\nabla^{2}f\succeq\mu\cdot I,\partial_{i}^{2}f\leq L_{i}$} & {\scriptsize Accelerated Coordinate Descent} & \multicolumn{2}{c|}{{\scriptsize Total Cost: $\sum_{i}\sqrt{\beta_{i}/\mu}\mathcal{T}_{\partial_{i}f}\log(\frac{1}{\epsilon})$}}\tabularnewline
\hline 
\end{tabular}
\par\end{centering}
\caption{Black-box Problems. The goal is to find $x$ such that $f(x)-\min_{x}f\protect\leq\epsilon$
with initial point $x^{(0)}=0$. $\mathcal{T}_{\nabla f}$ is the
cost of computing $\nabla f(x)$ for any $x$. $\mathcal{T}_{g}$
is the cost of computing $\arg\min_{x}g(x)+\frac{\lambda}{2}\|x-c\|^{2}$
for any $c$. $\mathcal{T}_{\Omega}$ is the cost of computing $\arg\min_{x\in\Omega}c^{\top}x$
for any $c$. $\mathcal{T}_{\nabla f_{\omega}}$, $\mathcal{T}_{\nabla f_{i}}$
and $\mathcal{T}_{\partial_{i}f}$ are the cost of computing $\nabla f_{\omega}$,
$\nabla f_{i}$ and $\partial_{i}f$ for one of the $\omega$ or $i$.\label{tab:Black-Box-Problems}}
\end{table}

\begin{table}[H]
\begin{centering}
\begin{tabular}{|>{\centering}m{0.1\paperwidth}|>{\centering}m{0.15\paperwidth}|>{\centering}m{0.15\paperwidth}|c|c|}
\hline 
{\scriptsize Convex Problem} & {\scriptsize Assumption} & {\scriptsize Algorithm} & {\scriptsize Depth} & {\scriptsize Total Work}\tabularnewline
\hline 
\hline 
\multirow{3}{0.1\paperwidth}{\centering{}{\scriptsize
\[
Ax=b
\]
}} & \centering{}{\scriptsize -} & {\scriptsize Fast Matrix Multiplication} & {\scriptsize$\log^{O(1)}(n)\log(\frac{1}{\epsilon})$} & {\scriptsize$n^{2.3728\cdots}\log(\frac{1}{\epsilon})$}\tabularnewline
\cline{2-5}
 & \centering{}{\scriptsize$\mu\cdot I\preceq A\preceq L\cdot I$} & {\scriptsize Conjugate Gradient} & {\scriptsize$\sqrt{\frac{L}{\mu}}\log(\frac{1}{\epsilon})$} & {\scriptsize$\nnz(A)\sqrt{\frac{L}{\mu}}\log(\frac{1}{\epsilon})$}\tabularnewline
\cline{2-5}
 & {\scriptsize$A_{ii}\geq\sum_{j\neq i}|A_{ji}|$} & {\scriptsize Sparse LU Factorization} & {\scriptsize$\log^{O(1)}(\frac{nL}{\mu})\log(\frac{1}{\epsilon})$} & {\scriptsize$\nnz(A)\log^{O(1)}(\frac{nL}{\mu})\log(\frac{1}{\epsilon})$}\tabularnewline
\hline 
\multirow{2}{0.1\paperwidth}{\centering{}{\scriptsize
\[
\min_{x}\|Ax-b\|_{2}
\]
}} & \centering{}{\scriptsize -} & {\scriptsize Subspace Embedding / Row Sampling} & {\scriptsize$\log^{O(1)}(n)\log(\frac{1}{\epsilon})$} & {\scriptsize$(\nnz(A)+r^{2.3728\cdots})\log(\frac{1}{\epsilon})$}\tabularnewline
\cline{2-5}
 & \centering{}{\scriptsize Each row of $A$ has $\leq2$ non-zeros} & {\scriptsize Sparse Cholesky Factorization} & {\scriptsize$\log^{O(1)}(n)\log(\frac{1}{\epsilon})$} & {\scriptsize$\nnz(A)\log^{O(1)}(n)\log(\frac{1}{\epsilon})$}\tabularnewline
\hline 
\multirow{2}{0.1\paperwidth}{\centering{}{\scriptsize
\[
\min_{Ax=b,x\in\Rn_{+}}c^{\top}x
\]
}} & \multirow{2}{0.15\paperwidth}{\centering{}{\scriptsize -}} & {\scriptsize Interior Point Method} & {\scriptsize$\sqrt{n}\log^{O(1)}(n)\log(\frac{1}{\epsilon})$} & {\scriptsize$n^{2.3728\cdots}\log(\frac{1}{\epsilon})$}\tabularnewline
\cline{3-5}
 &  & {\scriptsize Interior Point Method} & {\scriptsize$\sqrt{r}\log^{O(1)}(n)\log(\frac{1}{\epsilon})$} & {\scriptsize$\sqrt{r}(\nnz(A)+r^{2})\log^{O(1)}n\log(\frac{1}{\epsilon})$}\tabularnewline
\hline 
\centering{}{\scriptsize
\[
\min_{x}\sum_{i=1}^{n}f_{i}(a_{i}^{\top}x)
\]
} & \centering{}{\scriptsize -} & {\scriptsize Interior Point Method} & {\scriptsize$\sqrt{n}\log^{O(1)}(n)\log(\frac{1}{\epsilon})$} & {\scriptsize$n^{2.3728\cdots}\log(\frac{1}{\epsilon})$}\tabularnewline
\hline 
\end{tabular}
\par\end{centering}
\caption{Structured Problems. $r=\mathrm{rank}(A)$. We omit the precise definition
of $\epsilon$, but all algorithms have complexity depending on $\log(1/\epsilon)$.\label{tab:Structured-Problems}}
\end{table}


\section{Remarks on the Runtime}

\subsection{Structured vs black-box problems}

Table \ref{tab:Black-Box-Problems} shows the runtimes for some convex
problems assuming some restricted way to access the input convex function.
We call the way to access the function an \emph{oracle}. For example,
we say we only have a first-order oracle if we can only access the
function using the gradient of the function. In this setting, we are
particularly interested in how many times we need to use the oracle,
such as how many gradients we need to compute. Since the function
is only accessed through an oracle, we can usually bound how many
calls are needed to effectively ``learn'' the function. In fact,
all algorithms in the table are optimal in terms of the number of
oracle calls. This is in huge contrast to structured problems where
the tightness of all known results is open.

Table \ref{tab:Structured-Problems} has the runtimes for some general
convex problems. Unlike black-box problems, these algorithms assume
the full description of the problem. For this setting, we usually
aim for a runtime that depends on the structure of the problem instead
of the numerical values of the input. Usually, such algorithms converge
with a rate of $\log(1/\epsilon)$. For all of these problem, the
best possible runtime is open. Finally, we note that the last two
problems, linear programs and empirical risk minimization are complete
in the sense that they can be used to approximate any convex function.
Therefore, these last results seem to suggest that general convex
problems are as easy as general linear systems (for now!).

\subsection{Understanding the bounds}

All the runtimes should have the correct ``units''. For instance,
if the unit of the function value is $J$ (Joule) and the unit of
the space is $m$, then the units of $G$, $R$ and $\epsilon$ are
$Jm^{-1}$, $m$ and $J$ respectively. Therefore, the number of iterations
of gradient descent $\frac{GR}{\epsilon}$ is unitless. In the table,
we ignored some entries inside the $\log$ which leads to an ``incorrect''
unit. However, when everything is written correctly, the number of
iterations must be unitless. Furthermore, we can view $GR$, $LR^{2}$,
$\frac{G^{2}}{\mu}$ all as upper bounds of the initial error $f(x^{(0)})-f^{*}$.
For example, we know that
\[
f(x^{(0)})-f^{*}\leq\left\langle \nabla f(x^{(0)}),x^{(0)}-x^{*}\right\rangle \leq GR
\]
and that
\[
f(x^{(0)})-f^{*}\leq\left\langle \nabla f(x^{*}),x^{(0)}-x^{*}\right\rangle +\frac{L}{2}\|x^{(0)}-x^{*}\|^{2}\leq LR^{2}.
\]
Therefore, the terms $\frac{GR}{k}$, $\frac{G^{2}}{\mu k}$ and $\frac{LR^{2}}{k}$
are all just measuring how much we need to decrease multiplicatively
from the initial error. With this view, we can think the number of
iteration required for Lipschitz convex functions, strongly convex
functions and smooth functions are exactly $\frac{1}{\epsilon_{\mathrm{relative}}^{2}}$
, $\frac{1}{\epsilon_{\mathrm{relative}}}$ and $\frac{1}{\sqrt{\epsilon_{\mathrm{relative}}}}$.

\subsection{$L/\mu$ can be dimension dependent}

The ratio $L/\mu$ heavily depends on the application. For some easy
problems, it can be as small as $O(1)$. However, it can be very large
for many natural problems. Here we give a few examples of natural
convex functions with huge $\frac{L}{\mu}$. One very simple example
is 
\[
f(x)=\frac{1}{2}x_{1}^{2}+\frac{1}{2}\sum_{k=1}^{n-1}(x_{k}-x_{k+1})^{2}+\frac{1}{2}x_{n}^{2}
\]
This problem appears everywhere from physical systems, graph problems,
signal denoising to statistics. For this problem, we have $L=\Theta(1)$
and $\mu=\Theta(\frac{1}{n^{2}})$. This example shows that you can
have large $\frac{L}{\mu}$ even with small description coefficients
in the problem. Here is another example:
\[
f(x)=\norm{Ax-b}_{2}^{2}
\]
where $A$ is a random $n\times n$ matrix with independent Gaussian
entries. For this problem, we have $L=\Theta(n)$ and $\mu=\Theta(\frac{1}{n})$
with high probability. For an extreme example, the function $|x|$
has both $L=+\infty$ and $\mu=0$. 

\subsection{Accuracy of analysis}

For gradient descent and accelerated gradient descent, the convergence
rate is asymptotically accurate if we take $L\sim\lambda_{\max}(\nabla^{2}f(x^{*}))$
and $\mu\sim\lambda_{\min}(\nabla^{2}f(x^{*}))$ and $f\in\mathcal{C}^{2}$.
This is because any $\mathcal{C}^{2}$ function locally looks like
a quadratic around its minimum. We take gradient descent as an example.
Note that gradient descent is invariant under any rotation and translation
transformation. Therefore, we can assume $x^{*}=0$ and that $\nabla^{2}f(x^{*})$
is a diagonal matrix with diagonals $\lambda_{1},\lambda_{2},\cdots,\lambda_{n}$.
When $x$ is close to the optimal $x^{*}=0$, gradient descent $x^{\new}=x-h\nabla f(x)$
is similar to a coordinate-wise iteration $x_{i}^{\new}=(1-h\lambda_{i})x_{i}$.
This gives a precise description of the convergence rate.

For more complicated algorithms, the theory of convergence rate and
cost per iteration sometimes gives a bad estimate of the reality.
For example, the interior point method often takes $O(\log n)$ iterations
to converge despite the $\sqrt{n}$ worst-case bound on its convergence
rate; the simplex method is often the best of choice for solving linear
programs despite its exponential complexity bound.

The cost per iterations sometimes can be misleading also. Each iteration
of interior point methods takes $\nnz(A)+\mathrm{rank}(A)^{2}$ in
theory while it takes nearly linear time in practice for sparse matrices
due to the effectiveness of Cholesky decomposition. The $k^{th}$
iteration of cutting plane methods often takes $nk+k^{O(1)}$ if implemented
correctly despite its super-linear bound.

Finally, the error of floating point computation gradually starts
to have an effect when the target accuracy $\epsilon$ is small.

\section{Relations Between Different First-order Methods}

In the table \ref{tab:Black-Box-Problems}, we omitted many cases.
For example, what is the runtime if we have a stochastic function
$\min_{x}\E_{\omega}f_{\omega}(x)$ such that $\E_{\omega}f_{\omega}(x)$
is strongly convex? What is the runtime if we have a finite sum problem
$\frac{1}{n}\sum_{i}f_{i}(x)$ which is not strongly convex? In this
section, we show that one can reduce runtimes between different cases.
To highlight the idea, we will be informal and ignore polylogarithmic
terms. We note that some of the reduction ideas are not just useful
for proving theorems, but are practical tricks to speed up algorithms
and to avoid having to know parameters. See \cite{lin2015universal,frostig2015regularizing,o2015adaptive,allen2016optimal}
for more reduction tricks for first-order methods.

\subsection*{From Strongly Convex to Convex}

Suppose we have an algorithm $\mathcal{A}(f,x^{(0)})$ that inputs
a $\mu$-strongly convex function $f$ such that $\|\nabla f(x)\|_{2}\leq G$
for the domain of $f$ (such as a ball of radius $R$ around $x^{(0)}$)\footnote{Any strongly convex function on $\Rn$ cannot be uniformly Lipschitz
on $\Rn$. Therefore, we must assume the domain is bounded in some
way.} and outputs a point $x$ such that $f(x)\leq\min_{x}f(x)+\epsilon$
in $\frac{G^{2}}{\mu\epsilon}$ calls to $\nabla f$. Equivalently,
$f(x)\leq\min_{x}f(x)+\frac{G^{2}}{\mu k}$ in $k$ calls to $\nabla f$.
Now, we show how to recover the bound $(\frac{GR}{\epsilon})^{2}$
by the following algorithm:
\begin{itemize}
\item Output $\mathcal{A}(\widetilde{f},x^{(0)})$ where $\widetilde{f}(x)=f(x)+\frac{G}{R\sqrt{2k}}\|x\|_{2}^{2}.$
\end{itemize}
To understand the result, we let $f_{\sigma}(x)=f(x)+\frac{\sigma}{2}\|x\|_{2}^{2}$
and let $x$ be the output of $\mathcal{A}(f_{\sigma})$. Note that
$f_{\sigma}$ is $\sigma$ strongly convex. Hence, in $k$ iterations,
we have that
\[
f_{\sigma}(x)\leq\min_{x}f_{\sigma}(x)+\frac{G^{2}}{\sigma k}\leq f_{\sigma}(x^{*})+\frac{G^{2}}{\sigma k}=f(x^{*})+\frac{\sigma}{2}R^{2}+\frac{G^{2}}{\sigma k}
\]
where $x^{*}$ is the minimum of $f(x)$. To get the best runtime
for the fix $k$, we choose $\sigma$ such that 
\[
\frac{\sigma}{2}R^{2}=\frac{G^{2}}{\sigma k}
\]
which leads to $\sigma=\sqrt{\frac{2}{k}}\frac{G}{R}$ and hence the
total error is $\sqrt{2}\frac{GR}{\sqrt{k}}$ which is exactly what
we expect to get (up to constant).

\subsection*{From Convex to Strongly Convex}

Suppose we have an algorithm $\mathcal{A}(f,x^{(0)})$ that inputs
a Lipschitz function $f$ such that $\|\nabla f(x)\|_{2}\leq G$ and
that $\|x^{(0)}-x^{*}\|_{2}\leq R$ and outputs a point $x$ such
that $f(x)\leq\min_{x}f(x)+\epsilon$ in $(\frac{GR}{\epsilon})^{2}$
calls to $\nabla f$. Equivalently, $f(x)\leq\min_{x}f(x)+\frac{GR}{\sqrt{k}}$
in $k$ calls to $\nabla f$. Now, we show how to recover the result
$\widetilde{O}(\frac{G^{2}}{\mu\epsilon})$ by the following algorithm:
\begin{itemize}
\item For $i=1,2,\cdots,\widetilde{O}(1)$
\begin{itemize}
\item $x^{(i)}\leftarrow\mathcal{A}(f,x^{(i-1)})$ for $\frac{G^{2}}{\mu\epsilon}$
iterations
\end{itemize}
\end{itemize}
By the guarantee of $\mathcal{A}$, we have that 
\begin{equation}
f(x^{(i)})\leq\min_{x}f(x)+\frac{G\|x^{(i-1)}-x^{*}\|_{2}}{\sqrt{k}}=\min_{x}f(x)+\sqrt{\mu\epsilon}\|x^{(i-1)}-x^{*}\|_{2}.\label{eq:convex2strongly}
\end{equation}
On the other hand, the strong convexity of $f$ shows that $f(x^{(i)})-\min_{x}f(x)\geq\frac{\mu}{2}\cdot\|x^{(i)}-x^{*}\|_{2}^{2}$.
Therefore, we have that
\[
\|x^{(i)}-x^{*}\|_{2}^{2}\leq2\sqrt{\frac{\epsilon}{\mu}}\|x^{(i-1)}-x^{*}\|_{2}.
\]
Therefore, if $\|x^{(i-1)}-x^{*}\|_{2}\geq8\sqrt{\frac{\epsilon}{\mu}}$,
then the distance to OPT decreases by at least a factor of two. After
$\widetilde{O}(1)$ iterations, we have that $\|x^{(i-1)}-x^{*}\|_{2}=O\left(\sqrt{\frac{\epsilon}{\mu}}\right)$
and hence (\ref{eq:convex2strongly}) shows that $f(x^{(i)})\leq\min_{x}f(x)+\sqrt{\mu\epsilon}\cdot O\left(\sqrt{\frac{\epsilon}{\mu}}\right)=\min_{x}f(x)+O(\epsilon)$.
\begin{xca}
Given an algorithm that can minimize convex $f$ with $\epsilon$
additive error in $\sqrt{\frac{LR^{2}}{\epsilon}}$ iterations, show
how to use this algorithm to minimize $f$ again with $\epsilon$
additive error in $\widetilde{O}\left(\sqrt{\frac{L}{\mu}}\log\left(\frac{LR^{2}}{\epsilon}\right)\right)$
iterations.
\end{xca}


\subsection*{}\selectlanguage{english}%

